\section{Proof constructors for fixed bounded formulas}
\label{sec:p05-pa-bounded-proof-builders}

We continue with the selected tag-19 proof formula of
\cref{def:p05-selected-direct-pa-graph}. The constructions below turn the
numeral instances of a fixed bounded formula into annotated derivations.
The universal case includes an arithmetic exhaustion proof and actual
quantifier inferences. All functions have the finite-description meaning of
\cref{def:p05-checker-description-language}; they are verified in
\cref{sec:annex-p05-pa-bounded-proof-builders}.

\begin{hbdefinition}{def:p05-bounded-builder-conventions}{Fixed formulas, marked bounds and raw proof notation}
Fix a bounded formula in its marked presentation. Its primitive atomic
formulas are equalities, with bottom and the connectives
\(\neg,\land,\lor,\to\) also permitted. Any other propositional
abbreviation is expanded first. The construction additionally treats the
inherited comparisons as designated auxiliary literal blocks, in particular
in the marked guards below; their unbounded existential expansions are not
thereby reclassified as primitive atomic formulas of the hierarchy.
The marked quantifiers are \(\exists x_i\mathrel\diamond t\) and
\(\forall x_i\mathrel\diamond t\), with
\(\diamond\in\{<,\leq\}\) and \(x_i\) absent from the term \(t\).
Every comparison is expanded using the binder chosen in the original
expression. Every marked quantifier retains its displayed binder. Thus, for
a supplied comparison binder \(h\), put
\[
 G_<(i,h,u)=\mathsf{Lt}_h(v_i,u),\qquad
 G_{\leq}(i,h,u)=\mathsf{Le}_h(v_i,u).
\]
In an application \(h\ne i\), and the term in the third position has no
variable \(x_h\). The marked expansions have the exact codes
\[
 \mathsf{Ex}(i,G_\diamond(i,h,t)\mathbin{\&}c),\qquad
 \mathsf{All}(i,G_\diamond(i,h,t)\Rightarrow c).
\]
Here and below a term letter in a code expression denotes its term code.

Use \(\mathsf J,\mathsf L,\mathsf R,\mathsf{FE},\mathsf{FI},
\mathsf{EI},\mathsf{EE},\mathsf{II},\mathsf{NI},\mathsf B,
\mathsf{And},\mathsf{EEq},\mathsf{Sym},\mathsf{Tr}\), and
\(\mathsf{MP}^{\rm row}\) exactly as defined in
\cref{def:p05-builder-row-templates}. Write \(c\mathbin{|}d=N_9(c,d)\).
The following additional names are fixed row expressions:
\[
\begin{aligned}
 \mathsf{P}_0(R)&=\mathsf J_2(a_0(f(R)),0;R),&
 \mathsf{P}_1(R)&=\mathsf J_3(a_1(f(R)),0;R),\\
 \mathsf{OL}(R,d)&=\mathsf J_4(f(R)\mathbin{|}d,0;R),&
 \mathsf{OR}(c,R)&=\mathsf J_5(c\mathbin{|}f(R),0;R),\\
 \mathsf{OE}(H,R,S,j,k)&=\mathsf J_6(f(R),[j,k];H,R,S),&
 \mathsf{EF}(R,c)&=\mathsf J_{11}(c,0;R),\\
 \mathsf{RA}(R,c,k)&=\mathsf J_{12}(c,[k];R),&
 \mathsf{Keep}(H,R)&=\mathsf P_1(\mathsf{And}(H,R)).
\end{aligned}
\]
The projections \(\mathsf P_0,\mathsf P_1\) return zero unless the root
has tag 8 and arity two. The other expressions have their displayed values
on all numerical inputs; their successful-tree assertions require their
actual rule premises. In particular \(\mathsf{Keep}\) retains all open
leaves of both arguments, including a temporary assumption that must later
be discharged. No repetition or vacuous-discharge rule is added.

Write \(\mathsf{Safe}(R,c)\) for the arithmetic invariant (7) of
\cref{annex:p5.numeral-arithmetic-builders}: \(\mathsf{Tree}(R)=1\),
\(f(R)=c\), and every surviving leaf is unmarked and belongs to the
specific nine-entry list \(\mathcal A_*\). A root of a safe raw tree may
have free variables. Packing it is a selected proof when its root is a
sentence. In particular, on the sentence-and-axiom domain,
\[
 \mathsf{pack}(\mathsf{MP}^{\rm row}(R,S))
       =\mathsf{mp}_{19}(\mathsf{pack}(S),\mathsf{pack}(R)).
\]
The raw operation takes the antecedent first; the packed operation takes the
implication first.
\end{hbdefinition}

\begin{hbdefinition}{def:p05-bounded-em-template}{An explicit excluded-middle tree}
For a formula code \(c\), let \(d=c\mathbin{|}\mathsf{Neg}(c)\), and put
\[
\begin{aligned}
 T(c)&=\mathsf{NI}\bigl(
   \mathsf B(\mathsf{OL}(\mathsf L(c,1),\mathsf{Neg}(c)),
                    \mathsf L(\mathsf{Neg}(d),2)),c,0\bigr),\\
 \mathsf{EM}(c)&=\mathsf{RA}\bigl(
   \mathsf B(\mathsf{OR}(c,T(c)),\mathsf L(\mathsf{Neg}(d),2)),d,1\bigr).
\end{aligned}
\]
Every occurrence of a tree inserts a separate copy. Marker 1 denotes label
0, and marker 2 denotes label 1. These expressions also specify the values
on nonformulas, without a successful-tree assertion there.
\end{hbdefinition}

\begin{hblemma}{thm:p05-bounded-logical-templates}{proved=annex:p5.bounded-logical-templates}{Logical templates and nonvacuous use}
In \(\PA\), the new row expressions satisfy the exact local tests and
open-list updates in \cref{def:p05-checker-row-description} whenever their
displayed rule premises hold. For a formula \(c\), the tree
\(\mathsf{EM}(c)\) has root \(c\mathbin{|}\mathsf{Neg}(c)\) and no
open leaves. The two appearances of \(\mathsf{Neg}(d)\) in its last
premise are both discharged. The operation \(\mathsf{Keep}(H,R)\)
has root \(f(R)\), with precisely the shifted leaves of \(H\) followed
by those of \(R\).
\end{hblemma}

\begin{hbdefinition}{def:p05-bounded-sum-exhaustion}{Finite exhaustion of a numeral sum}
Define the formula-code recursion, with counter last,
\[
 \mathcal D(t,0)=\mathsf{bc},\qquad
 \mathcal D(t,Sn)=\mathcal D(t,n)\mathbin{|}\mathsf{Eq}(t,\nu(n)).
\]
Thus zero cases are represented by bottom, and an additional case is appended
on the right. Define two separate orientations, indexed here by the fixed
metalinguistic choice \(\epsilon\in\{0,1\}\):
\[
\begin{array}{c|c|c|c}
 \epsilon&g_\epsilon(t,u)&Z_\epsilon(t)&S_\epsilon(t,u)\\\hline
 0&a(t,u)&\mathsf A_4(t)&\mathsf A_5(t,u)\\
 1&a(u,t)&\mathsf A_Z(t)&\mathsf A_L(u,t).
\end{array}
\]
The last two columns are the actual instance chains and copied derived
records from \cref{def:p05-builder-axiom-instances}. They have roots
\(g_\epsilon(t,o)=t\) and
\(g_\epsilon(t,s(u))=s(g_\epsilon(t,u))\), respectively.

We define \(\mathsf C_\epsilon(x,z,w,n)\) for pairwise distinct
\(x,z,w\geq3\). Its failure value on any other triple is zero. The
following stage expression \(F_\epsilon(x,z,w,n,R)\) is a fixed
composition with a zero/successor branch on \(n\). Set
\[
\begin{aligned}
 e&=\mathsf{Eq}(v_z,o),&
 a_n&=\mathsf{Eq}(g_\epsilon(v_x,v_z),\nu(n)),&
 d_n&=\mathcal D(v_x,Sn),\\
 H&=\mathsf L(a_n,1),& Z&=\mathsf L(e,2),&
 K&=\mathsf L(\mathsf{Eq}(v_z,s(v_w)),4).
\end{aligned}
\]
The zero alternative for \(x_z\) is the tree
\[
 U=\mathsf{OR}\left(\mathcal D(v_x,n),
   \mathsf{Tr}\left(
    \mathsf{Tr}(\mathsf{Sym}(Z_\epsilon(v_x)),
      \mathsf{Sym}(\mathsf{Cg}_{u\mapsto g_\epsilon(v_x,u)}(Z))),H\right)\right).
\]
In the predecessor alternative form
\[
 W=\mathsf{Tr}\left(\mathsf{Sym}(S_\epsilon(v_x,v_w)),
       \mathsf{Tr}(\mathsf{Sym}
           (\mathsf{Cg}_{u\mapsto g_\epsilon(v_x,u)}(K)),H)\right).
\]
Its intended root is
\(\mathsf{Eq}(s(g_\epsilon(v_x,v_w)),\nu(n))\). Set
\[
 V=\begin{cases}
  \mathsf{EF}(\mathsf B(W,\mathsf A_1(g_\epsilon(v_x,v_w))),d_n),&n=0,\\
  \mathsf{OL}\left(
    \mathsf{MP}^{\rm row}\left(
      \mathsf{MP}^{\rm row}(W,
          \mathsf A_2(g_\epsilon(v_x,v_w),\nu(D(n)))),
      \mathsf{FE}(\mathsf{FE}(R,v_x),v_w)\right),
                           \mathsf{Eq}(v_x,\nu(n))\right),&n>0.
 \end{cases}
\]
Define
\[
\begin{aligned}
 V'&=\mathsf{EE}\bigl(
    \mathsf{MP}^{\rm row}(\mathsf L(\mathsf{Neg}(e),3),\mathsf A_3(v_z)),
                 V,1,\mathsf{Eq}(v_z,s(v_1)),w,3\bigr),\\
 F_\epsilon(x,z,w,n,R)&=
  \mathsf{FI}\bigl(\mathsf{FI}\bigl(
   \mathsf{II}(\mathsf{OE}(\mathsf{EM}(e),U,V',1,2),a_n,0),z\bigr),x\bigr).
\end{aligned}
\]
Finally the promised recursion is
\[
 \mathsf C_\epsilon(x,z,w,0)=F_\epsilon(x,z,w,0,0),\qquad
 \mathsf C_\epsilon(x,z,w,Sn)
       =F_\epsilon(x,z,w,Sn,\mathsf C_\epsilon(x,z,w,n)).
\]
Include the triple's validity guard at the base and step; the guard is
constant in the recursion counter. These are two finite descriptions, not a
recursion over orientations or over formulas.
\end{hbdefinition}

\begin{hblemma}{thm:p05-bounded-sum-exhaustion}{proved=annex:p5.bounded-sum-exhaustion}{Internal exhaustion invariant in both orientations}
For each orientation, \(\PA\) proves, for all pairwise distinct
\(x,z,w\geq3\) and all \(n\),
\[
 \mathsf{Safe}\left(\mathsf C_\epsilon(x,z,w,n),
  \mathsf{All}\left(x,\mathsf{All}\left(z,
    \mathsf{Eq}(g_\epsilon(v_x,v_z),\nu(n))
                       \Rightarrow\mathcal D(v_x,Sn)\right)\right)\right).
\]
In particular these are selected sentence proofs after packing. The
surviving axiom leaves are precisely of the already allowed kinds in
\(\mathcal A_*\); no additional induction-axiom family is used by the
output trees. The verification uses induction in \(\PA\) on the
numerical recursion counter.
\end{hblemma}

\begin{hbdefinition}{def:p05-bounded-range-exhaustion}{Strict sums and comparison exhaustion}
First define \(\mathsf{CS}(x,z,w,n)\), on the same triple domain, by
the following composition. Let
\[
 c_0=\mathsf{Eq}(a(v_x,s(v_z)),\nu(n)),\quad H=\mathsf L(c_0,1),\quad
 W=\mathsf{Tr}(\mathsf{Sym}(\mathsf A_5(v_x,v_z)),H).
\]
At \(n=0\) put \(V=\mathsf B(W,\mathsf A_1(a(v_x,v_z)))\).
At \(n>0\) put
\[
 V=\mathsf{MP}^{\rm row}\left(
    \mathsf{MP}^{\rm row}(W,\mathsf A_2(a(v_x,v_z),\nu(D(n)))),
    \mathsf{FE}(\mathsf{FE}(\mathsf C_0(x,z,w,D(n)),v_x),v_z)\right).
\]
Return \(\mathsf{FI}(\mathsf{FI}(\mathsf{II}(V,c_0,0),z),x)\).
Invalid triples return zero.

For \(i\ne h\), put \(M=\max(2,i,h)\) and
\(y=M+1,x=M+2,z=M+3,w=M+4\). For either fixed comparison
orientation let \(G=G_\diamond(i,h,\nu(n))=\mathsf{Ex}(h,c)\),
\(H=\mathsf L(G,1)\), and \(K=\mathsf L(\mathsf{sub}(c,h,v_y),2)\).
Use
\[
 Q=\begin{cases}
   \mathsf C_1(x,z,w,n),&\diamond=\leq,\\
   \mathsf{CS}(x,z,w,n),&\diamond=<,
 \end{cases}
\qquad
 V=\mathsf{MP}^{\rm row}(K,\mathsf{FE}(\mathsf{FE}(Q,v_i),v_y)).
\]
Define
\[
 \mathsf{E}_\diamond(i,h,n)
       =\mathsf{II}(\mathsf{EE}(H,V,h,c,y,1),G,0).
\]
This is a separate composition for each \(\diamond\); set its value to
zero when \(i=h\). Set \(N_<(n)=n\), \(N_{\leq}(n)=Sn\).
\end{hbdefinition}

\begin{hblemma}{thm:p05-bounded-range-exhaustion}{proved=annex:p5.bounded-range-exhaustion}{Actual exhaustion of a marked range}
In \(\PA\), on the stated index domains, \(\mathsf{CS}\) is safe with
root
\[
 \mathsf{All}(x,\mathsf{All}(z,
    \mathsf{Eq}(a(v_x,s(v_z)),\nu(n))\Rightarrow\mathcal D(v_x,n))),
\]
and \(\mathsf E_\diamond(i,h,n)\) is safe with root
\[
 G_\diamond(i,h,\nu(n))\Rightarrow
                        \mathcal D(v_i,N_\diamond(n)).
\]
The strict zero case derives bottom from the assumed comparison. The weak
zero case retains the single possible equality to zero. Every existential
elimination satisfies its all-variable test and discharges a surviving
minor assumption with its specified formula and marker.
\end{hblemma}

\begin{hbdefinition}{def:p05-bounded-case-elimination}{The finite case tree and the quantified constructor}
Let \(c\) be a formula code, let \(i\) be a variable index, and let
\(Q\) be a list of raw proofs. Define a recursion with counter last:
\[
\begin{aligned}
 \mathsf{Cases}(c,i,Q,0)&=\mathsf{EF}(\mathsf L(\mathsf{bc},1),c),\\
 T_n&=\mathsf{EEq}\bigl(
       \mathsf{Sym}(\mathsf L(\mathsf{Eq}(v_i,\nu(n)),1)),e(Q,n),
                                      \nu(n),v_i,c,i\bigr),\\
 \mathsf{Cases}(c,i,Q,Sn)&=
   \mathsf{OE}\bigl(\mathsf L(\mathcal D(v_i,Sn),1),
                            \mathsf{Cases}(c,i,Q,n),T_n,0,0\bigr),\\
 \mathsf{CI}(c,i,Q,n)&=
   \mathsf{II}(\mathsf{Cases}(c,i,Q,n),\mathcal D(v_i,n),0).
\end{aligned}
\]
The entry \(T_n\) is the step expression, with the preceding state supplied
only to the other minor premise. These expressions are total; successful
use requires \(\mathsf{Safe}(e(Q,k),\mathsf{sub}(c,i,\nu(k)))\) for
every \(k<n\).

Now supply a variable-free term code \(u\), a numerical value \(b\),
and a raw tree \(R_u\) safe with root \(\mathsf{Eq}(u,\nu(b))\).
Let \(i\ne h\), \(G=G_\diamond(i,h,u)\),
\(N=N_\diamond(b)\), and \(j=1+\max(i,h)\). Put
\[
\begin{aligned}
 T&=\mathsf{EEq}(R_u,\mathsf L(G,1),u,\nu(b),G_\diamond(i,h,v_j),j),\\
 V&=\mathsf{MP}^{\rm row}
       (T,\mathsf E_\diamond(i,h,b)),\\
 W&=\mathsf{MP}^{\rm row}(V,\mathsf{CI}(c,i,Q,N)),\\
 \mathsf U_\diamond(c,i,h,u,b,R_u,Q)&=
                  \mathsf{FI}(\mathsf{II}(W,G,0),i).
\end{aligned}
\]
Thus the bound is normalized within the temporary guard; the final
implication still has its original guard code \(G\).
\end{hbdefinition}

\begin{hbtheorem}{thm:p05-bounded-case-elimination}{proved=annex:p5.bounded-case-elimination}{From all numeral cases to one quantified derivation}
In \(\PA\), if \(c\) is a formula, \(\ell(Q)\geq n\), and every
\(k<n\) entry of \(Q\) is safe with root
\(\mathsf{sub}(c,i,\nu(k))\), then
\(\mathsf{CI}(c,i,Q,n)\) is safe with root
\(\mathcal D(v_i,n)\Rightarrow c\).

Suppose additionally that \(c\) has no free indices other than \(i\),
that \(u\) is variable-free, that \(i\ne h\), and that \(R_u\)
has the stated safe equality root. If the case premise holds through
\(N_\diamond(b)\), then \(\mathsf U_\diamond\) is safe with exact root
\[
             \mathsf{All}(i,G_\diamond(i,h,u)\Rightarrow c).
\]
This is a sentence, and its packed record satisfies
\(\mathsf{Prf}^{19}_{\PA}\). The assertions include \(n=0\),
\(b=0\), unused free indices, bound-variable shadowing in \(c\), and
the originally supplied comparison binder \(h\).
\end{hbtheorem}

\begin{hbdefinition}{def:p05-fixed-bounded-builders}{The two proof descriptions for each fixed bounded formula}
For a fixed formula \(\varphi\), list its distinct free indices in
increasing order. Numerical inputs \(\bar m\) supply their values.
Write \(\varphi^\nu(\bar m)\) for its literal numeral substitution
code, and \(b_\varphi(\bar m)\) for the following fixed zero-one
description. At equality and comparison use the numerical term values and
the tests of \cref{def:p05-numeral-builder-notation}; bottom has value
zero. Negation uses \(Z\), conjunction uses product, disjunction uses
sign of the sum, and implication uses
\(I(b_\psi,b_\chi,1)\) at \(\psi\to\chi\).
At a marked quantifier with bound value \(b=t^\#(\bar m)\), let
\(N=N_\diamond(b)\). Existential quantification uses
\(\mathsf{sg}(\sum_{k<N}b_\psi(\bar m,k))\), and universal
quantification uses \(\prod_{k<N}b_\psi(\bar m,k)\).

At a child call, its free-index list receives the parent's numerical values
at the corresponding indices; the currently bound index receives \(k\).
It overrides that index only within this child. This is a fixed tuple
projection at each node of the one fixed formula. For a quantifier, the
bound is evaluated before its fold. A nested bound can consequently depend
on all its surrounding numerical witnesses. No input ranges over formula
descriptions.

Define raw descriptions \(R_\varphi^+,R_\varphi^-\) in the same fixed
formation, returning zero outside their respective success conditions
\(b_\varphi=1\) and \(b_\varphi=0\). For equality and comparisons
take the raw record inside the corresponding packed literal description in
\cref{def:p05-fixed-term-literal-builders}; decode the outer tags 19 then
15, returning zero on a failed pattern. This retains the exact six literal
roots and their supplied comparison binders. Bottom uses zero for the
positive description and \(R_{\neg\bot}\) for the negative description.

Here are the complete connective clauses. In them \(c=\psi^\nu\),
\(d=\chi^\nu\); each child proof has the unchanged tuple arguments,
and \(H\) means \(\mathsf L(\varphi^\nu,1)\) in a negative clause.
Write \(H_c=\mathsf L(c,1)\) and \(H_d=\mathsf L(d,1)\).
\begin{enumerate}
\item At \(\neg\psi\), the positive tree is \(R_\psi^-\), and the
negative tree is
\[
 \mathsf{NI}(\mathsf B(R_\psi^+,\mathsf L(\mathsf{Neg}(c),1)),
                                         \mathsf{Neg}(c),0).
\]
\item At \(\psi\land\chi\), the positive tree is
\(\mathsf{And}(R_\psi^+,R_\chi^+)\). The negative tree is
\[
 \begin{cases}
 \mathsf{NI}(\mathsf B(\mathsf P_0(H),R_\psi^-),c\mathbin{\&}d,0),&b_\psi=0,\\
 \mathsf{NI}(\mathsf B(\mathsf P_1(H),R_\chi^-),c\mathbin{\&}d,0),&b_\psi=1.
 \end{cases}
\]
\item At \(\psi\lor\chi\), the positive tree is
\(\mathsf{OL}(R_\psi^+,d)\) when \(b_\psi=1\), and
\(\mathsf{OR}(c,R_\chi^+)\) when \(b_\psi=0\). For the negative
tree put
\[
 K_\psi=\mathsf B(H_c,R_\psi^-),\qquad
 K_\chi=\mathsf B(H_d,R_\chi^-),
\]
and return
\[
 \mathsf{NI}(\mathsf{OE}(H,K_\psi,K_\chi,0,0),c\mathbin{|}d,0).
\]
\item At \(\psi\to\chi\), the positive tree is
\[
 \begin{cases}
 \mathsf{II}(\mathsf{EF}(\mathsf B(H_c,R_\psi^-),d),c,0),&b_\psi=0,\\
 \mathsf{II}(\mathsf{Keep}(H_c,R_\chi^+),c,0),&b_\psi=1.
 \end{cases}
\]
The negative tree is
\[
 \mathsf{NI}\bigl(
   \mathsf B(\mathsf{MP}^{\rm row}(R_\psi^+,H),R_\chi^-),c\Rightarrow d,0\bigr).
\]
\end{enumerate}
The outer success guard precedes each listed subcase test. Thus the second
negative-conjunction subcase has \(b_\chi=0\); the second
positive-implication subcase has \(b_\chi=1\).

For a marked quantifier let \(u=t^\nu(\bar m)\),
\(b=t^\#(\bar m)\), \(N=N_\diamond(b)\), and let \(c\) be the
matrix code after replacing its other free parameters by their numerals,
leaving the current index \(i\) free. Set
\(G=G_\diamond(i,h,u)\), \(C=G\mathbin{\&}c\), and
\(A=\mathsf{All}(i,G\Rightarrow c)\). The bound term's proof is
\(R_t(\bar m)\) from \cref{def:p05-fixed-term-literal-builders}.
For every \(k<N\), let \(L_k\) be the raw positive fixed-literal
proof of \(\mathsf{sub}(G,i,\nu(k))\). Its free parameter tuple
includes the current \(k\) before numeral substitution; its comparison
binder is the original \(h\). This is the already fixed literal
description applied through tuple projections, not a new proof search.

For the positive existential, choose the least \(k<N\) with
\(b_\psi(\bar m,k)=1\), with failure value zero, and use
\[
 \mathsf{EI}(\mathsf{And}(L_k,R_\psi^+(\bar m,k)),i,C,\nu(k)).
\]
For the negative universal, choose the least \(k<N\) with
\(b_\psi(\bar m,k)=0\), with failure value zero, and use
\[
 \mathsf{NI}\left(\mathsf B\left(
    \mathsf{MP}^{\rm row}(L_k,\mathsf{FE}(\mathsf L(A,1),\nu(k))),
                          R_\psi^-(\bar m,k)\right),A,0\right).
\]
Both searches are the earlier bounded least-search description and are
called only in their outer success branches.

Let \(Q^+\), respectively \(Q^-\), be the length-\(N\) fold starting
with zero and appending \(R_\psi^+(\bar m,k)\), respectively
\(R_\psi^-(\bar m,k)\), at step \(k\). For the positive universal use
\[
               \mathsf U_\diamond(c,i,h,u,b,R_t(\bar m),Q^+).
\]
For the negative existential first put
\[
 U=\mathsf U_\diamond(\mathsf{Neg}(c),i,h,u,b,R_t(\bar m),Q^-),
 \quad E=\mathsf{Ex}(i,C),\quad y=1+\max(2,m(E)),
 \quad K=\mathsf L(\mathsf{sub}(C,i,v_y),2).
\]
Use the tree
\[
 \mathsf{NI}\left(\mathsf{EE}\left(\mathsf L(E,1),
    \mathsf B\left(\mathsf P_1(K),
       \mathsf{MP}^{\rm row}(\mathsf P_0(K),\mathsf{FE}(U,v_y))\right),
                                    i,C,y,1\right),E,0\right).
\]
The two displayed uses of \(K\) are two copies; the minor discharge
removes both. Finally define
\(d_\varphi^\pm(\bar m)=I(b,\mathsf{pack}(R_\varphi^\pm(\bar m)),0)\),
where \(b=b_\varphi\) for the plus sign and \(b=Z(b_\varphi)\)
for the minus sign. These guards also preserve a zero failure value after
packing.
\end{hbdefinition}

\begin{hbtheorem}{thm:p05-fixed-bounded-builders}{proved=annex:p5.fixed-bounded-builders}{Internal proof construction for each fixed bounded formula}
For each fixed bounded formula \(\varphi\), \(\PA\) proves
\[
\begin{aligned}
 &b_\varphi(\bar m)\in\{0,1\},\qquad
             b_\varphi(\bar m)=1\ \leftrightarrow\ \varphi(\bar m),\\
 &b_\varphi(\bar m)=1\ \longrightarrow\
             \mathsf{Safe}(R_\varphi^+(\bar m),\varphi^\nu(\bar m)),\\
 &b_\varphi(\bar m)=0\ \longrightarrow\
             \mathsf{Safe}(R_\varphi^-(\bar m),\mathsf{Neg}(\varphi^\nu(\bar m))),
\end{aligned}
\]
with universal numerical closure. Here \(\varphi(\bar m)\) is the
original fixed arithmetic formula at those numerical variables. Consequently
it proves both
\[
\begin{aligned}
 \forall\bar m\bigl(\varphi(\bar m)\to\exists p\,
  (\Gamma_{d_\varphi^+}(\bar m,p)\land
                \mathsf{Prf}^{19}_{\PA}(p,\varphi^\nu(\bar m)))\bigr),\\
 \forall\bar m\bigl(\neg\varphi(\bar m)\to\exists p\,
  (\Gamma_{d_\varphi^-}(\bar m,p)\land
       \mathsf{Prf}^{19}_{\PA}(p,\mathsf{Neg}(\varphi^\nu(\bar m))))\bigr).
\end{aligned}
\]
All graph applications in these assertions expand to their specified
\(\Gamma\) formulas. The construction preserves literal roots, all
original binders, dependent bounds and witnesses, and nonvacuous discharge.
An empty free-parameter list gives a nullary description and an empty
quantifier prefix. This theorem supplies the bounded-formula dependency;
the subsequent evidence-to-proof construction must still specialize it and
construct the original unbounded existential and other evidence clauses
before an assertion of D3 for the selected predicate is available.
\end{hbtheorem}
